\chapter{Perfil Profissional do Egresso}
\label{chp:perfil_egresso}

A formação do(a) \ppctitulacao envolve uma base técnico-científica em eletricidade, eletrônica, instrumentação, controle e automação, com ênfase em sistemas embarcados, redes industriais, acionamentos elétricos e, como diferencial deste curso, robótica e inteligência artificial aplicada à automação (\ppcenfasecurricular). Além da competência técnica, o curso promove o desenvolvimento de habilidades interpessoais, de liderança técnica, comunicação, trabalho em equipe e pensamento crítico, associadas a uma consciência ética, social e ambiental.

Este capítulo descreve exclusivamente o perfil profissional do egresso do \ppccurso. Embora o \ppccurso integre a ABI com o Curso Superior de Tecnologia em Automação Industrial, e os cinco primeiros períodos comuns às duas graduações desenvolvam competências e conteúdos de fundamentação compartilhados (\autoref{sec:integracao_pedagogica_abi}), cada curso mantém perfil de egresso próprio: o perfil descrito a seguir corresponde integralmente ao \ppcgrau em \ppcenfasecurricular conferido por este curso, e não se confunde com o perfil profissional do(a) Tecnólogo(a) em Automação Industrial, definido no respectivo PPC. O estudante que concluir as duas graduações (\autoref{chp:abi}) deverá atender, separadamente, ao perfil e aos requisitos formativos de cada uma delas.

Em consonância com as Diretrizes Curriculares Nacionais do Curso de Graduação em Engenharia~\cite{cne_ces_2_2019,cne_ces_1_2021}, espera-se que o egresso deste curso:

\begin{itemize}
    \item Atue com autonomia técnica, iniciativa e compromisso com a atualização profissional contínua, dada a rápida evolução das tecnologias de automação;
    \item Reconheça e se comprometa com a segurança do trabalho, a sustentabilidade e a eficiência energética em suas práticas profissionais;
    \item Compreenda os impactos sociais, éticos e ambientais da automação industrial nos processos produtivos e no mundo do trabalho;
    \item Domine os fundamentos técnicos para especificar, projetar, implementar, operar e manter sistemas de automação industrial robustos, seguros e eficientes;
    \item Comunique-se com clareza e eficiência com públicos técnicos e não técnicos;
    \item Seja capaz de atuar de forma responsável e proativa em contextos colaborativos e interdisciplinares, em equipes de manutenção, engenharia e produção;
    \item Exerça liderança técnica na condução de projetos de automação e na orientação de equipes de operação e manutenção.
\end{itemize}

São competências esperadas do egresso do \ppccurso, conforme definidas na aba \texttt{Competencias} da matriz curricular deste perfil:

\begin{itemize}
    \item Projetar, especificar e sintonizar sistemas de controle analógico e digital para processos industriais;
    \item Programar controladores lógicos programáveis (CLPs) e sistemas supervisórios (SCADA/IHM) para automação de processos;
    \item Especificar, instalar e calibrar sensores, transdutores e atuadores utilizados em instrumentação industrial;
    \item Projetar e dimensionar sistemas de acionamento elétrico de máquinas e processos industriais;
    \item Integrar dispositivos e sistemas por meio de redes de comunicação industrial e protocolos de automação;
    \item Aplicar normas de segurança do trabalho e de sistemas instrumentados de segurança em ambientes industriais automatizados;
    \item Atuar profissionalmente com responsabilidade ética, social e ambiental;
    \item Desenvolver soluções com sistemas embarcados aplicados ao controle e à automação de processos (competência não obrigatória, associada ao diferencial em sistemas inteligentes do curso);
    \item Planejar e gerir projetos de automação industrial, incluindo aspectos de custo, prazo e qualidade (competência não obrigatória).
\end{itemize}

O título profissional conferido é o de \textbf{\ppctitulacao}, com exercício profissional regulamentado pelo Sistema CONFEA/CREA nos termos da Lei nº 5.194/1966~\cite{brasil_lei_5194_1966} (\autoref{sec:exercicio_profissional}). A Resolução CONFEA nº 1.073/2016\footnote{Texto oficial: \url{https://normativos.confea.org.br/Ementas/Visualizar?id=59111}.} estabelece a sistemática geral de atribuição de títulos, atividades e competências com base na formação, e a Resolução CONFEA nº 1.156/2025~\cite{confea_1156_2025} consolida as atribuições da modalidade Eletricista. O egresso deste curso estará habilitado para atuar nas mais diversas áreas da automação industrial, como:

\begin{itemize}
    \item Instrumentação, metrologia e controle de processos industriais;
    \item Automação industrial — controladores lógicos programáveis (CLP), interfaces homem-máquina (IHM) e sistemas supervisórios (SCADA);
    \item Redes industriais, protocolos de automação e Internet Industrial das Coisas (IIoT);
    \item Acionamentos e máquinas elétricas, hidráulica e pneumática industrial;
    \item Robótica industrial e manufatura automatizada;
    \item Sistemas embarcados aplicados a dispositivos ciberfísicos de automação;
    \item Inteligência artificial e análise de dados aplicadas a processos industriais;
    \item Cibersegurança de redes e sistemas de controle industrial;
    \item Gestão, manutenção e qualidade de processos industriais automatizados.
\end{itemize}

O perfil do egresso reflete o compromisso do curso com uma formação orientada por competências, fundamentada em valores de excelência técnica, responsabilidade ética, segurança e transformação da prática industrial.

\section{Egresso do Curso de \ppccurso}

O egresso do \ppccurso da Universidade Federal de Uberlândia é um(a) profissional com sólida formação técnico-aplicada, preparado(a) para atuar no projeto, na implementação, na integração, na operação e na manutenção de sistemas automatizados industriais. Sua formação o(a) capacita a unir eletrônica, controle, redes industriais e sistemas inteligentes de forma crítica, aplicada, ética e responsável, com capacidade de adaptação a diferentes plantas e processos produtivos.

A formação ofertada permite ao egresso desenvolver competências integradas nos seguintes eixos:

\begin{enumerate}
\item Fundamentos Técnico-Científicos

\begin{itemize}
    \item Dominar os conhecimentos fundamentais de matemática, física, eletricidade e eletrônica, integrando teoria e prática para o projeto e a análise de sistemas de automação;
    \item Especificar, implementar, testar, colocar em operação e manter sistemas de instrumentação, controle, acionamento e redes industriais, conforme boas práticas da engenharia aplicada;
    \item Aplicar o raciocínio lógico e a modelagem de sistemas dinâmicos para resolver problemas de controle e automação com eficiência, qualidade e segurança;
    \item Analisar a viabilidade técnica e econômica de soluções de automação, considerando requisitos de desempenho, custo, escalabilidade e sustentabilidade energética.
\end{itemize}

\item Inovação, Robótica e Inteligência Artificial Aplicada

\begin{itemize}
    \item Aplicar técnicas de robótica industrial, sistemas inteligentes e inteligência artificial na modernização de processos produtivos;
    \item Desenvolver e integrar soluções tecnológicas inovadoras em automação, unindo criatividade com conhecimento técnico-aplicado;
    \item Atuar de forma proativa em ambientes empreendedores e de inovação tecnológica, reconhecendo oportunidades de modernização de plantas industriais;
    \item Aprofundar-se, por meio da integralização de, no mínimo, duas Ênfases Formativas — Máquinas Inteligentes, Acionamentos e Proteção Industrial (MIAPI); Robótica Autônoma e Sistemas de Controle (RASC); e Sistemas Embarcados, IoT e Computação Inteligente (SEICI) —, em percursos técnicos coerentes diretamente relacionados a este eixo (\autoref{chp:modulos_optativos}).
\end{itemize}

\item Ética, Sustentabilidade e Impacto Social

\begin{itemize}
    \item Refletir criticamente sobre o impacto social, ambiental, ético e econômico da automação industrial nos processos produtivos e no mundo do trabalho;
    \item Projetar soluções que respeitem a segurança do trabalho, a eficiência energética e o uso racional de recursos;
    \item Agir com responsabilidade, autonomia técnica e consciência social, contribuindo para o desenvolvimento industrial sustentável.
\end{itemize}

\item Comunicação, Colaboração e Liderança Técnica

\begin{itemize}
    \item Comunicar-se com clareza, precisão e adequação com diferentes públicos, de forma oral e escrita, incluindo documentação técnica de projetos e processos;
    \item Trabalhar eficazmente em equipes multidisciplinares de engenharia, operação e manutenção, assumindo papéis de liderança técnica;
    \item Demonstrar responsabilidade compartilhada e respeito à diversidade no trabalho em equipe e na mediação de conflitos.
\end{itemize}

\item Aprendizagem Contínua e Adaptabilidade

\begin{itemize}
    \item Demonstrar autonomia na gestão do próprio aprendizado, reconhecendo a necessidade de atualização constante diante da rápida evolução das tecnologias de automação;
    \item Adaptar-se a novas plantas, ferramentas, protocolos e ambientes de trabalho industrial;
    \item Avaliar criticamente tecnologias emergentes de automação e aplicar princípios de inovação de forma ética e eficaz.
\end{itemize}

\item Atuação Profissional

\begin{itemize}
    \item Atuar de forma competente, ética e inovadora nos seguintes campos:

    \begin{itemize}
        \item Instrumentação e controle de processos industriais
        \item Automação industrial (CLP, IHM e SCADA)
        \item Redes industriais e Internet Industrial das Coisas (IIoT)
        \item Acionamentos elétricos, hidráulica e pneumática
        \item Robótica industrial e manufatura automatizada
        \item Inteligência artificial e ciência de dados aplicadas a processos industriais
        \item Cibersegurança de sistemas de controle industrial
    \end{itemize}

    \item Intervir com competência técnica em plantas e processos industriais locais, nacionais e, quando aplicável, internacionais, respeitando legislações, normas técnicas e expectativas sociais.
\end{itemize}
\end{enumerate}

\section{Habilidades e Competências}

O profissional egresso do \ppccurso da UFU deve desenvolver um conjunto articulado de conhecimentos técnico-científicos, habilidades práticas e disposições profissionais, de modo a atuar com excelência, responsabilidade social e capacidade de inovação, conforme estabelecem as Diretrizes Curriculares Nacionais do Curso de Graduação em Engenharia~\cite{cne_ces_2_2019,cne_ces_1_2021}. A referência anterior às Diretrizes da Educação Profissional e Tecnológica e ao CNCST, herdada do Curso Superior de Tecnologia, foi removida por não se aplicar a este Bacharelado.

\begin{enumerate}[label=\Roman*.]
\item Competências Técnicas

\textcolor{red}{[A LISTA DE COMPETÊNCIAS TÉCNICAS REPRODUZIDA NO INÍCIO DESTE CAPÍTULO E RETOMADA AQUI FOI HERDADA INTEGRALMENTE DO CURSO SUPERIOR DE TECNOLOGIA EM AUTOMAÇÃO INDUSTRIAL E DESCREVE COMPETÊNCIAS DE NÍVEL TECNÓLOGO (OPERAÇÃO, PROGRAMAÇÃO E ESPECIFICAÇÃO DE SISTEMAS JÁ EXISTENTES) — PENDENTE DE SUBSTITUIÇÃO PELO NDE/COLEGIADO POR COMPETÊNCIAS TÉCNICAS PRÓPRIAS DE UM BACHARELADO EM ENGENHARIA (PROJETO, MODELAGEM, ANÁLISE E DIMENSIONAMENTO DE SISTEMAS ORIGINAIS), CONFORME AS DIRETRIZES CURRICULARES NACIONAIS DE ENGENHARIA, ANTES DA SUBMISSÃO DESTE PPC]}

O curso visa assegurar ao egresso as competências técnicas listadas na aba \texttt{Competencias} da matriz curricular (reproduzidas no início deste capítulo), com destaque para o projeto e a sintonia de sistemas de controle, a programação de CLPs e sistemas supervisórios, a especificação de instrumentação industrial, o projeto de acionamentos elétricos e a integração de redes industriais.

\item Habilidades Cognitivas e Práticas

Complementando as competências técnicas, espera-se que o egresso demonstre:

\begin{itemize}
    \item Raciocínio lógico, técnico e analítico para resolver problemas de automação e controle;
    \item Capacidade de leitura e interpretação de diagramas elétricos, fluxogramas de processo e documentação técnica;
    \item Comunicação eficaz com públicos técnicos e não técnicos;
    \item Capacidade de colaboração em equipes multidisciplinares de engenharia, operação e manutenção;
    \item Adaptabilidade a novos contextos tecnológicos e industriais;
    \item Organização e gestão do tempo e do próprio aprendizado;
    \item Liderança técnica em projetos de automação, com tomada de decisão embasada e eficaz;
    \item Capacidade de avaliação crítica de sistemas automatizados e de propostas de modernização industrial.
\end{itemize}

\item Disposições e Atitudes Profissionais

Além dos conhecimentos e habilidades, busca-se que o egresso adote atitudes compatíveis com a atuação ética e transformadora, tais como:

\begin{itemize}
    \item Comprometimento com a aprendizagem contínua e o aprimoramento técnico-profissional;
    \item Responsabilidade social, ambiental e ética no uso e na implementação de tecnologias de automação;
    \item Respeito à diversidade, à equidade e à inclusão nos ambientes de trabalho;
    \item Rigor técnico, atenção à segurança do trabalho e cultura de prevenção de riscos industriais;
    \item Proatividade, colaboração e escuta ativa em equipes de trabalho;
    \item Engajamento com o impacto positivo da automação na produtividade e na segurança dos processos industriais.
\end{itemize}
\end{enumerate}

Essa integração entre os referenciais normativos garante que o egresso do \ppccurso da UFU esteja preparado não apenas para os desafios técnicos da profissão, mas também para atuar com responsabilidade social e liderança técnica em nível local e regional.

\section{Exercício Profissional no Brasil}
\label{sec:exercicio_profissional}

O exercício, no Brasil, das profissões de Engenheiro, Arquiteto e Engenheiro-Agrônomo é regulamentado desde 1966~\cite{brasil_lei_5194_1966}. A verificação e a fiscalização do exercício profissional nas áreas por ela abrangidas são exercidas pelo Conselho Federal de Engenharia e Agronomia (CONFEA) e pelos Conselhos Regionais de Engenharia e Agronomia (CREA). Os profissionais habilitados na forma estabelecida nessa legislação só poderão exercer a profissão após o registro no Conselho Regional sob cuja jurisdição se achar o local de sua atividade.

Tanto o(a) Tecnólogo(a) em Automação Industrial, formado(a) pelo Curso Superior de Tecnologia em Automação Industrial, quanto o(a) \ppctitulacao, formado(a) pelo \ppccurso, podem estar vinculados ao Sistema CONFEA/CREA. Pertencer ao mesmo sistema profissional, todavia, não significa possuir o mesmo título ou as mesmas atribuições: cada curso confere um título profissional distinto, e as atribuições correspondentes são analisadas pelo Sistema CONFEA/CREA conforme a formação efetivamente demonstrada — para a qual são relevantes o PPC, a matriz curricular, as ementas, as cargas horárias e as atividades práticas de cada curso, além do perfil do egresso aqui descrito. O egresso deste curso poderá requerer registro no Sistema CONFEA/CREA, observados o título profissional correspondente, a regulamentação vigente e a análise da formação efetivamente desenvolvida no curso — não se promete, neste PPC, atribuição profissional automática ou idêntica à do(a) Tecnólogo(a) em Automação Industrial.

As atribuições profissionais dos(as) egressos(as) deste curso, na modalidade Eletricista do Sistema CONFEA/CREA — que abrange, entre outras, as áreas de automação, computação, software e produção industrial —, são atualmente consolidadas pela Resolução CONFEA nº 1.156, de 24 de outubro de 2025~\cite{confea_1156_2025}, que reorganizou em um único instrumento normativo as atribuições antes fragmentadas nos arts. 8º e 9º da Resolução CONFEA nº 218/1973~\cite{confea_218_1973} e nas Resoluções CONFEA nº 380/1993 e nº 427/1999 (revogadas). A Resolução CONFEA nº 1.073/2016, indicada na fonte oficial acima, aplica-se à sistemática de atribuições profissionais também para esta Engenharia; a Resolução nº 313/1986 é a norma específica dos Tecnólogos. O exercício profissional está sujeito, ainda, ao Código de Ética Profissional adotado pela Resolução CONFEA nº 1.002/2002~\cite{confea_1002_2002}. A extensão concreta das atribuições será definida pelo CREA a partir da formação efetivamente comprovada; este PPC não assegura atribuições automáticas para robótica ou inteligência artificial.

A ocupação da Classificação Brasileira de Ocupações (CBO) mais próxima da denominação deste curso é a de código 2021-10, ``Engenheiro de controle e automação'' (sinônimos: engenheiro de automação, engenheiro de controle, engenheiro de instrumentação), distinta das ocupações de Tecnólogo — 2021-15 (Tecnólogo em Mecatrônica) e 2021-20 (Tecnólogo em Automação), próprias do Curso Superior de Tecnologia em Automação Industrial e listadas no Catálogo Nacional de Cursos Superiores de Tecnologia, que não se aplicam a este Bacharelado em Engenharia. \textcolor{red}{[CONFIRMAR JUNTO AO NDE/COLEGIADO SE O CÓDIGO CBO 2021-10 É O MAIS ADEQUADO PARA A ÊNFASE EM ROBÓTICA E INTELIGÊNCIA ARTIFICIAL DESTE CURSO, OU SE HÁ CÓDIGO MAIS ESPECÍFICO A SER ADOTADO]}
